%% kapa: exam=1 round=33 year=2022 period=1 form=A subject=real_estate subject_ko=부동산학원론 seq=25 number=105
%% source: Q-Net 감정평가사 기출문제 (cst00309) artlSeq=5209997 file=raw/1008/5209997/2022년도 제33회 감정평가사 1차시험 1교시(민법,경제학,부동산학).pdf
%% conversion: kordoc 4.24.51 -> tools/gen1.py (auto); review=proofread-sonnet; flags=none
\begin{question}{105}
부동산개발에 관한 설명으로 옳은 것을 모두 고른 것은?
\begin{framedtext}
ㄱ. 부동산개발업의 관리 및 육성에 관한 법률상 부동산개발은 토지를 건설공사의 수행 또는 형질변경의 방법으로 조성하는 행위 및 건축물을 건축, 대수선, 리모델링 또는 용도를 변경하거나 공작물을 설치하는 행위를 말하며, 시공을 담당하는 행위는 제외한다. \par
ㄴ. 혼합방식은 개발전의 면적·등급·지목 등을 고려하여, 개발된 토지를 토지소유주에게 종전의 토지위치에 재분배하는 것을 말한다. \par
ㄷ. 흡수율분석은 수요·공급분석을 통하여 대상부동산이 언제 얼마만큼 시장에서 매각 또는 임대될 수 있는지를 파악하는 것이다. \par
ㄹ. 개발권양도제(TDR)는 일정하게 주어진 개발허용한도 내에서 해당 지역의 토지이용규제로 인해 사용하지 못하는 부분을 다른 지역에 양도할 수 있는 것이다.
\end{framedtext}
\begin{choices}
  \item ㄱ, ㄷ
  \item ㄷ, ㄹ
  \item ㄱ, ㄴ, ㄹ
  \item ㄱ, ㄷ, ㄹ
  \item ㄴ, ㄷ, ㄹ
\end{choices}
\end{question}
\answer{4}
